\chapter{Concentration et loi des grands nombres}
\section*{Exercice 23.1 -- Bienaymé--Tchebychev}
[Base] Une variable a moyenne 50 et variance 9. Majorer $P(|X-50|\ge6)$.

\section*{Exercice 23.2 -- Taille d’échantillon}
[Méthode] Pour une fréquence de Bernoulli indépendantes de paramètre inconnu, garantir $P(|F_n-p|\ge0{,}05)\le0{,}01$ par Tchebychev.

\section*{Exercice 23.3 -- Rôle de l’indépendance}
[Critique] Si $X_1=\cdots=X_n=X$ avec $X$ Bernoulli de paramètre $1/2$, la moyenne se concentre-t-elle autour de $1/2$ ?

\section*{Exercice 23.4 -- Convergence en probabilité}
[Preuve] Si $E(X_n)=m$ et $V(X_n)\le C/n^2$, démontrer $X_n\to m$ en probabilité.

\section*{Exercice 23.5 -- Intervalle aléatoire}
[Raisonnement] Pour une fréquence $F_n$, justifier que l'intervalle aléatoire $[F_n-\varepsilon,F_n+\varepsilon]$ contient $p$ avec probabilité au moins $1-1/(4n\varepsilon^2)$.

\section*{Exercice 23.6 -- Attente géométrique}
[Approfondissement] Des essais indépendants de succès $p\in]0,1[$ sont répétés. Pour le rang $T$ du premier succès, calculer $P(T>n)$ et $P(T>m+n\mid T>m)$.
